\documentclass[10pt,a4paper]{amsart}
\usepackage[margin=19mm]{geometry}
\usepackage{amsmath,amssymb,mathtools}
\usepackage[scr=rsfs,cal=euler]{mathalfa}
\usepackage{xfrac}
\usepackage[unicode,hidelinks,bookmarks=false]{hyperref}
\newcommand{\ie}{i.e.\ }
\newcommand{\cf}{cf.\ }
\newcommand{\resp}{respectively\ }
\newcommand{\Subsection}{\subsection}
\providecommand{\nobreakdash}{\nobreak\hbox{-}}
\setlength{\emergencystretch}{2em}
\multlinegap0pt
\usepackage{tikz}
\usepackage{amssymb}
\usepackage{float}
\usepackage{xcolor}
\usepackage[shortlabels]{enumitem}
\usepackage{enumerate}
\newtheorem{lemma}{Lemma}
\def\B{\mathcal{T}}
\def\Lint{\mathsf{Lint}}
\def\Sint{\mathsf{Sint}}
\def\l{\mathsf{left}}
\def\leaf{\mathsf{leaf}}
\def\lint{\mathsf{lint}}
\def\maj{\mathsf{maj}}
\def\sint{\mathsf{sint}}
\title{The $q$-analogues of $\gamma$-positivity of Eulerian polynomials via group actions}
\author{Taifeng Ding, Lintong Wang, Sherry H.F. Yan}
\date{Source: arXiv:2609.12405v1 (CC BY 4.0)}
\begin{document}
\maketitle

\noindent\emph{Source and licence.} The $q$-analogues of $\gamma$-positivity of Eulerian polynomials via group actions. Version arXiv:2609.12405v1 (CC BY 4.0), distributed by arXiv under the Creative Commons Attribution 4.0 International licence, http://creativecommons.org/licenses/by/4.0/, as recorded in the arXiv licence metadata for this exact version and shown on the abstract page https://arxiv.org/abs/2609.12405v1. Full licence notice: \texttt{licenses/arxiv-2609.12405-CC-BY-4.0.txt}.\par\smallskip

\noindent\textit{Editorial scope.} Counted unit: the article's lemma lem-psi-1 (source0.tex lines 1863--1878). The article's complete original proof of it is delivered with it (source0.tex lines 1879--1985, one \texttt{proof} environment of 107 lines). No mathematics is written, repaired or re-derived: the statement and the proof are the article's own lines, in the article's own notation, and no retained line is substituted or re-flowed.  Retained uncounted as the source-local context the proof needs: lines 1278--1288; lines 1389--1392.  Nothing was omitted from any retained span, so the package carries no omission record.  No retained line contains a citation, so the wrapper carries no bibliography and no bibliography entry.  No retained line contains a figure environment, an image-inclusion command or drawing code, so no image is delivered and the package declares no asset dependency.  Editorial formatting only, in addition: an AMS \texttt{amsart} preamble that restates the article's own declarations and macros used by the retained lines, namely the environments \texttt{lemma} on separate counters, together with the standard packages those lines need.  The retained environments print in this document as Lemma 1 in order of appearance, whereas the article prints the same items with the numbering of its own full document.  The complete native source of the article as distributed in the arXiv e-print for this exact version is bundled unchanged as \texttt{sources/210113/source0.tex}.
	 	\begin{lemma}\label{lem-tree}
	 		Given a binary tree $T\in \B_n$, we have
	 		\begin{equation}\label{eq-sint}
	 			\sint(T)=n+1-2\leaf(T),
	 		\end{equation}
	 		and 
	 		\begin{equation}\label{eq-left}
	 			\l(T)= \lint(T)+ \leaf(T)-1,
	 		\end{equation}
	 		where  $\sint(T)$  denotes the number of singleton internal nodes    of $T$.
	 	\end{lemma}

	 	\begin{lemma}\label{lem-alpha}
	 		Let $T\in \B_n$. For any $x\in [n]$, we have 
	 		$\alpha(T)=\alpha(\theta_x(T))$. 
	 	\end{lemma}

	 	\begin{lemma}\label{lem-psi-1}
	 		
	 	Fix $n\geq 2$. 	For  any tree $T\in \B_{n-1}$ and $0\le c\le n-1$, we have
	 		\[
	 		\l(\psi(T, c))=\left\{ 
	 		\begin{array}{ll}
	 			\l(T) & \mbox{if  $0\leq c\leq \l(T)$},\\[2pt]
	 			\l(T)+1 & \mbox{otherwise},
	 		\end{array}
	 		\right.
	 		\]
	 		and
	 		\[
	 		\maj(\psi(T, c))=\maj(T)+c.
	 		\]
	 		\end{lemma}

	 		\begin{proof}
	 			Here we retain the notations in the definition of the map $\psi$. Let $Q=\psi(T,c)$.   We proceed  by considering  four cases.
	 			
	 			\noindent{\bf Case 1: $0\leq c\leq k-1.$ }\\
	 			By Lemma \ref{lem-tree}, we have $\l(T)=\lint(T)+\leaf(T)-1$.  Then $c\leq k-1=\leaf(T)-1$ would imply that $c\leq \l(T)$.  Now we proceed to show that 
	 			$$
	 			\maj(Q)=\maj(T)+c \,\, \mbox{and}\,\, \l(Q)=\l(T).
	 			$$
	 		 
	 			Recall that $P$ is obtained from $T$ by connecting the node $\ell_{c+1}$ and node $n$ by a right edge.  It is plain to see that $$\l(T)=\l(P),$$ $$\leaf(T)=\leaf(P),$$   and $$\alpha(P)=\alpha(T)+c.$$  
	 			Recall that $Q$ is obtained from $P$ by adjusting the child-type (left/right)   of each singleton internal node of $P$.  Moreover, we have  $\beta(Q)=(\beta_1, \beta_2, \ldots, \beta_{n-2k}, 0)$ and $\beta(T)=(\beta_1, \beta_2, \ldots, \beta_{n-2k})$.  By Lemma \ref{lem-alpha},   adjusting the child-type (left/right)   of each singleton internal node of $P$ in the procedure from $P$ to $Q$ does not affect the statistic $\alpha(P)$.   It is easily seen that the procedure from $P$ to $Q$ preserves the number of leaves and the number of left edges. This gives that  $$\alpha(Q)=\alpha(P)=\alpha(T)+c,$$
	 			 $$\leaf(Q)=\leaf(P)=\leaf(T),$$
	 			 and $$\l(Q)=\l(P)=\l(T).
	 			 $$
	 			  Hence, we deduce that
	 			 $$
	 			 \begin{array}{lll}
	 			 \maj(Q)&=&\alpha(Q)+\sum\limits_{i=1}^{n-2k+1}(\leaf(Q)+i-1)\beta_i\\
	 			 &=&\alpha(T)+c+\sum\limits_{i=1}^{n-2k}(\leaf(T)+i-1)\beta_i\\
	 			 &=& \maj(T)+c,
	 			 \end{array}
	 			 $$
	 			as desired.
	 			
	 		 
	 		 	\noindent{\bf Case 2: $n-k\leq c\leq n-1.$ }\\
	 		 By Lemma \ref{lem-tree}, we have $\l(T)=\lint(T)+\leaf(T)-1\leq n-2k+k-1=n-k-1$.  Then $c\geq n-k$ would imply that $c> \l(T)$.  Now we proceed to show that 
	 		 $$
	 		 \maj(Q)=\maj(T)+c \,\, \mbox{and}\,\, \l(Q)=\l(T)+1.
	 		 $$
	 		 
	 		 Recall that $P$ is obtained from $T$ by connecting the node $\ell_r$ and node $n$ by a left edge where $r=c-n+k+1$.  It is plain to see that $$\l(P)=\l(T)+1,$$ $$\leaf(T)=\leaf(P),$$  and  $$\alpha(P)=\alpha(T)+r-1.$$  
	 		 Recall that $Q$ is obtained from $P$ by adjusting the child-type (left/right)   of each singleton internal node of $P$.  Moreover, we have  $\beta(Q)=(\beta_1, \beta_2, \ldots, \beta_{n-2k}, 1)$ and $\beta(T)=(\beta_1, \beta_2, \ldots, \beta_{n-2k})$.  By Lemma \ref{lem-alpha},   adjusting the child-type (left/right)   of each singleton internal node of $P$ in the procedure from $P$ to $Q$ does not affect the statistic $\alpha(P)$.   It is easily seen that the procedure from $P$ to $Q$ preserves the number of leaves and the number of left edges. This gives that  $$\alpha(Q)=\alpha(P)=\alpha(T)+r-1,$$
	 		 $$\leaf(Q)=\leaf(P)=\leaf(T)=k,$$
	 		 and $$\l(Q)=\l(P)=\l(T)+1.
	 		 $$
	 		 Hence, we deduce that
	 		 $$
	 		 \begin{array}{lll}
	 		 	\maj(Q)&=&\alpha(Q)+(\leaf(Q)+n-2k)+\sum\limits_{i=1}^{n-2k}(\leaf(Q)+i-1)\beta_i\\
	 		 	&=&\alpha(T)+r-1+k+n-2k+\sum\limits_{i=1}^{n-2k}(\leaf(T)+i-1)\beta_i\\
	 		 	&=& \maj(T)+c
	 		 \end{array}
	 		 $$
	 		 as desired.
	 		 
	 		 	\noindent{\bf Case 3: $k\leq c\leq k+s-1.$ }\\
	 		 By Lemma \ref{lem-tree}, we have $\l(T)=\lint(T)+\leaf(T)-1=k+s-1$.  Then $c\leq k+s-1$ would imply that $c\leq  \l(T)$.  Now we proceed to show that 
	 		 $$
	 		 \maj(Q)=\maj(T)+c \,\, \mbox{and}\,\, \l(Q)=\l(T).
	 		 $$
	 		 Recall that $Q$ is obtained from $T$ by connecting the node $u_{i_r}$ and node $n$ by a right edge where $r=c-k+1$.  It is plain to see that $$\l(Q)=\l(T),$$ 
	 		 $$\leaf(Q)=\leaf(T)+1$$ and   $$\alpha(Q)=\alpha(T)+2\leaf(T)+i_r-1.$$    Moreover, we have  $$\beta(Q)=(\beta_1, \beta_2, \ldots, \beta_{i_r-1},   \beta_{i_r+1}, \ldots,  \beta_{n-2k})$$ and $$\beta(T)=(\beta_1, \beta_2, \ldots, \beta_{i_r-1}, 1,  \beta_{i_r+1}, \ldots,  \beta_{n-2k}).$$  
	 		 Hence, we deduce that
	 		 $$
	 		 \begin{array}{lll}
	 		 	\maj(Q)&=&\alpha(Q)+\sum\limits_{j=1}^{i_r-1}(\leaf(Q)+j-1)\beta_j+\sum\limits_{j=i_r+1}^{n-2k}(\leaf(Q)+j-2)\beta_j\\
	 		 	&=&\alpha(T) +2\leaf(T)+i_r-1+\sum\limits_{j=1}^{i_r-1}(\leaf(T)+j)\beta_j+\sum\limits_{j=i_r+1}^{n-2k}(\leaf(T)-1+j)\beta_j\\
	 		 	&=&  \alpha(T) +\leaf(T)+\sum\limits_{j=1}^{i_r-1}\beta_j+\sum\limits_{j=1}^{n-2k}(\leaf(T)+j-1)\beta_j\\
	 		  
	 		 	&=&  \alpha(T) +k+r-1+\sum\limits_{j=1}^{n-2k}(\leaf(T)+j-1)\beta_j\\
	 		 	&=& \maj(T)+c 
	 		 \end{array}
	 		 $$
	 		 as desired.
	 		 
	 			\noindent{\bf Case 4: $k+s\leq c\leq n-k-1.$ }\\
	 		By Lemma \ref{lem-tree}, we have $\l(T)=\lint(T)+\leaf(T)-1=k+s-1$.  Then $c\geq k+s$ would imply that $c>  \l(T)$.  Now we proceed to show that 
	 		$$
	 		\maj(Q)=\maj(T)+c \,\, \mbox{and}\,\, \l(Q)=\l(T)+1.
	 		$$
	 		
	 		
	 		
	 		Recall that $Q$ is obtained from $T$ by connecting the node $u_{j_r}$ and node $n$ by a left edge where $r=c-k-s+1$.   Clearly, we have
	 		\begin{equation}\label{eq-proof-1}
	 			\# \{u_i\mid i<j_r, u_i\in \Lint(T)\}+\# \{u_i\mid i>j_r, u_i\in \Lint(T)\}=s, 
	 		\end{equation}
	 		\begin{equation}\label{eq-proof-2}
	 			\# \{u_i\mid i>j_r, u_i\in \Lint(T)\}+r= \# \{u_i\mid i\geq j_r, u_i\in \Sint(T)\}=n-2k-j_r+1,
	 		\end{equation}
	 		and 
	 		\begin{equation}\label{eq-proof-3}
	 			\# \{u_i\mid i<j_r, u_i\in \Lint(T)\}=\sum\limits_{i=1}^{j_r-1}\beta_i.
	 		\end{equation}
	 		Invoking (\ref{eq-proof-1})-(\ref{eq-proof-3}), we deduce that
	 		\begin{equation}\label{eq-proof-4}
	 			n-2k-j_r+\sum\limits_{i=1}^{j_r-1}\beta_i=r+s-1.
	 		\end{equation}
	 		
	 		  It is plain to see that $$\l(Q)=\l(T)+1,$$ 
	 		$$\leaf(Q)=\leaf(T)+1$$ and   $$\alpha(Q)=\alpha(T)+\leaf(T)+n-2k-j_r.$$    Moreover, we have  $$\beta(Q)=(\beta_1, \beta_2, \ldots, \beta_{j_r-1},   \beta_{j_r+1},  \ldots, \beta_{n-2k})$$ and $$\beta(T)=(\beta_1, \beta_2, \ldots, \beta_{j_r-1}, 0,  \beta_{j_r+1}, \ldots,  \beta_{n-2k}).$$  
	 		Hence, we deduce that
	 		$$
	 		\begin{array}{lll}
	 			\maj(Q)&=&\alpha(Q)+\sum\limits_{i=1}^{j_r-1}(\leaf(Q)+i-1)\beta_i+\sum\limits_{i=j_r+1}^{n-2k}(\leaf(Q)+i-2)\beta_i\\
	 			&=&\alpha(T)+\leaf(T) +n-2k-j_r+\sum\limits_{i=1}^{j_r-1}(\leaf(T)+i)\beta_i+\sum\limits_{i=j_r+1}^{n-2k}(\leaf(T)-1+i)\beta_i\\
	 			&=&  \alpha(T)+\leaf(T) +n-2k-j_r+\sum\limits_{i=1}^{j_r-1}\beta_i+\sum\limits_{i=1}^{n-2k}(\leaf(T)+i-1)\beta_i\\
	 			
	 			&=&  \alpha(T) +k+r+s-1+\sum\limits_{i=1}^{n-2k}(\leaf(T)+i-1)\beta_i\\
	 			&=& \maj(T)+c 
	 		\end{array}
	 		$$
	 		as desired, where the fourth equality follows from \eqref{eq-proof-4}.  This completes the proof.
	 		 
	 		 
	 			\end{proof}

\end{document}
